\section{The sunflower bound}\label{sec:intro}

How many sets can a family contain before some \(r\) of them meet in
exactly the same core? Such a configuration is a \emph{sunflower}.
More precisely, distinct sets \(A_1,\ldots,A_r\) form an \(r\)-sunflower
if \(A_i\cap A_j=C\) for every \(i\ne j\). Removing the core \(C\)
leaves pairwise disjoint petals. The core can be empty, as it is
for any \(r\) disjoint sets.

Erd\H{o}s and Rado asked whether bounding the number of elements in
each set forces a sunflower once the family is exponentially large.
The ground set may be arbitrarily large. Write \(M_r(k)\) for the
largest size of a family of distinct \(k\)-element sets containing
no \(r\)-sunflower.

\begin{theorem}\label{thm:main}
For each integer \(r\ge3\) there is a constant \(B_r\) such that every
family of distinct sets of size at most \(k\) containing no
\(r\)-sunflower has at most \(B_r^k\) members, for every integer
\(k\ge1\). In particular \(M_r(k)\le B_r^k\). With logarithms to
base two, one may take
\[
 B_r=2^{C H^4 L^{5/2}},\qquad H=\log_2(2r),\quad L=\log_2(2H),
\]
where \(C\) is an absolute constant.
\end{theorem}

The classical argument gives \(M_r(k)\le k!(r-1)^k\)~\cite{ER}.
Take a maximal collection of disjoint members. It has at most \(r-1\)
sets, and its union, of size at most \(k(r-1)\), meets every member.
Among the sets containing a particular element, deleting that element
gives a sunflower-free family of \((k-1)\)-sets. Summing over the union
and iterating gives the factorial. Theorem~\ref{thm:main} gives a base
independent of \(k\).

Exponential growth is necessary even on a ground set of size \(k(r-1)\).
Partition it into \(k\) blocks of size \(r-1\) and take every set that
chooses one element from each block. There are \((r-1)^k\) such sets.
In a sunflower, the choices within a block would have to be all equal
or all distinct. With \(r\) choices and only \(r-1\) elements they would
all be equal, in every block, giving the same set \(r\) times. Thus the
optimal base \(b_r=\sup_{k\ge1}M_r(k)^{1/k}\) is at least \(r-1\).
Whether \(b_r=O(r)\) remains open. The formulation for sets of size at
most \(k\) follows from the uniform one by padding each set to size
\(k\) with new elements private to that set. Padding preserves all
pairwise intersections.

\subsection{The obstruction in one coordinate}\label{sec:overview}

Consider first a family in which every set chooses one element from
each of \(k\) disjoint classes. List those elements in class order
to obtain a word of length \(k\). In any \(r\) such words, a column
with a common symbol contributes an element to the core. If its symbols
are all different, its elements belong to separate petals. The
obstruction is a \emph{partial repetition}: a symbol occurs at least
twice but not in all \(r\) words. Figure~\ref{fig:columns} shows why a
single such column prevents a sunflower. A random coloring of the
ground set reduces the general \(k\)-uniform problem to these families
at a factor of at most \(k^k/k!\le e^k\) in their size
(Section~\ref{sec:words}). It suffices to consider finite subfamilies,
since any violation of the bound already occurs in one.

\begin{figure}[!htb]
\centering
\begin{tikzpicture}[x=.87cm,y=.62cm,font=\small]
  \foreach \offset/\third/\highlight/\heading in {0/f/0/{A sunflower},5.8/d/1/{After one entry changes}}{
    \begin{scope}[xshift=\offset cm]
      \node[font=\small\bfseries] at (2.0,3.75) {\heading};
      \foreach \x in {1,2,3}{\node[text=ink] at (\x,2.96) {\(\x\)};}
      \draw[rounded corners=2pt,draw=sea,fill=sea!6]
                         (.58,-.41) rectangle (1.42,2.41);
      \ifnum\highlight=1
        \draw[rounded corners=2pt,draw=rust,fill=rust!8]
                          (1.58,-.41) rectangle (2.42,1.41);
      \fi
      \foreach \y/\which/\xone/\xtwo/\xthree in
        {2/1/a/b/c,1/2/a/d/e,0/3/a/\third/g}{
        \node[anchor=east] at (.36,\y) {\(A_{\which}\)};
        \node at (1,\y) {\(\xone\)};
        \node at (2,\y) {\(\xtwo\)};
        \node at (3,\y) {\(\xthree\)};
      }
    \end{scope}
  }
\end{tikzpicture}
\caption{In the left family every pair of sets meets in \(\{a\}\).
Changing \(f\) to \(d\) in the last row makes \(A_2\cap A_3=\{a,d\}\).
The other two intersections still equal \(\{a\}\).}\label{fig:columns}
\end{figure}

Call a column \emph{admissible} if it is constant or has all symbols
distinct. In a sunflower-free family, the only ordered \(r\)-tuples
admissible in every column consist of a single repeated word. Indeed,
a column of distinct symbols would make the words distinct, hence give
a sunflower. If every column is constant, the words agree.

The proof turns this observation into a count. Give each of \(r\)
groups of sites an independently chosen word from a family of size \(N\).
At each coordinate apply a tensor operation which kills every partial
repetition and returns a fixed nonzero tensor on a constant column.
These two requirements suffice. Every tuple other than a single repeated word
has a partial repetition somewhere, and the \(N\) repeated-word tuples
all survive. The operations have converted \(r\) independently summed
word labels into a single label shared across the groups.

To measure what this conversion costs, we use a numerical function on
tensors whose value is multiplicative under tensor products and can only
decrease under local linear maps. The logarithm of its value is a
\emph{price}. Sharing a label with \(N\) possible values among the sites
of \(P\) has price \(w(P)\log_2 N\). The Euclidean geometry of the
asymptotic spectrum~\cite{oaG} lets us choose a single price for which
the rates \(w(P)\) are normalized mean widths of the corresponding
sets of points. Place the groups so that the rate for their union
exceeds the sum of their separate rates by \(\rho>0\). The change from
independent labels to a shared one then costs at least \(\rho\log_2 N\).
If the net cost of testing a column is at most \(f\), we obtain
\[
                          \rho\log_2 N\le kf .
\]
The central difficulty is to keep \(f\) bounded for an alphabet of
\emph{any} size. Even two symbols being compared may have arbitrarily
long binary names.

\par\Needspace{7\baselineskip}
\subsection{From a fixed petal count to the quantitative bound}

Part~\ref{part:fixed} proves the fixed-petal bound in the plane. If any
two symbols agree, all remaining symbols must agree with them for the
column to be admissible. Choosing two groups as anchors reduces this
requirement to a list of implications. When the two anchors agree,
they can copy their common symbol directly to a site between them.
The other implications use an equality test that shortens two binary
labels by pairing the alphabet in all possible ways of the form \(\{x,x+p\}\). Two different labels become equal
under exactly one of these pairings. The proof reconstructs and
subtracts that erroneous term. In the planar placement, the net costs
of all the corrections sum to a finite bound even as the label length
grows (Theorem~\ref{thm:filter}). This planar construction already
gives \(\log B_r=O(r^5)\).

Part~\ref{part:quant} obtains the stated dependence on \(r\) by treating
a whole column at once. Its test again shortens each symbol, preserving
every genuine repetition. A partially repeated column can now pass
only by becoming entirely constant. The task is to cancel exactly
those false agreements. The shortened labels are obtained by removing leading bits and
shifting the remaining bits by a mask. If a repeated symbol lies in a
small class with the same leading bits, a recursive test on just that
class pays for the correction. When almost all the active sites outside
the anchors (the satellites) carry the same symbol, splitting those sites into two separated
classes gives a majority in both. The vote in each class can be sent
to a designated site at a cost per output bit governed by one wire
(Lemma~\ref{lem:vote}). The Euclidean placement gives the two votes
together enough width to test their common value against the anchors.
This gives a bound uniform in the alphabet, with \(\log B_r\) bounded
by a polynomial in \(\log r\).
Corollary~\ref{cor:weighted} extends the count to word families that
contain sunflowers. The cost of excluding them is controlled by the rank
of an auxiliary tensor that vanishes on the sunflower tuples and equals
one on the repeated-word tuples.

\subsection{Earlier bounds and related sunflower problems}\label{sec:history}

Erd\H{o}s and Rado suggested replacing the factorial in their 1960
bound by \(c^k\) for an absolute \(c\)~\cite[p.~86]{ER}. In the
notation of Theorem~\ref{thm:main}, this would give \(b_r=O(r)\).
Erd\H{o}s later stated the fixed-petal conjecture explicitly, allowing
the base to depend arbitrarily on \(r\)~\cite[Section~4]{ERstrong}.
For three petals the block-transversal lower bound can be improved.
A construction of Abbott, Hanson and Sauer gives
\(b_3\ge\sqrt{10}\)~\cite{AHS}.

After earlier improvements by Kostochka~\cite{Kostochka}, Alweiss,
Lovett, Wu and Zhang brought the base down to a polylogarithmic
function of \(k\) for fixed \(r\)~\cite{ALWZ}. Their argument finds a
subfamily whose members have a common intersection \(C\), such that
a random subset of the ground set is likely to contain one of the
nonempty remainders \(A\setminus C\). If a random color class in
an \(r\)-coloring contains no remainder with probability less than
\(1/r\), some coloring puts one remainder in each class. These
\(r\) disjoint remainders give the petals. Their examples also show why a logarithmic loss
arises here. For fixed sampling and success probabilities there are
families of size \((\log k)^{(1-o(1))k}\) with no such robust
sunflower~\cite[Section~3]{ALWZ}. Rao's coding argument gave
\(M_r(k)\le (Cr\log(rk))^k\)~\cite{Rao}. Bell, Chueluecha and Warnke
refined the probabilistic step to give \((Cr\log(k+1))^k\)~\cite{BCW}.
Related random-containment estimates enter the study of thresholds of
increasing families~\cite{Thresholds,PP}. Sunflower lemmas also play
a role in Razborov's lower bounds for monotone
circuits~\cite{Razborov}. Rao surveys these connections~\cite{RaoSurvey}.

For \(k\)-uniform families of VC dimension \(d\ge1\), Ge, Wang, Xu and Zhao
obtained the width-independent base \(50dr\)~\cite[Theorem~1.1]{VC},
building on earlier results for restricted set
systems~\cite{FPS,vcsmall}. Shangguan improved the base to
\(O(r\log(d+1))\)~\cite[Corollary~1.4]{Shangguan}. For unrestricted
families Fukuyama proposed a base \(O(r^4)\)~\cite{Fukuyama}.
A product-family obstruction to the sequence-selection estimate in
that argument is described in~\cite{Plasma}.

The size \(N\) of the ground set gives a different measure of the
problem. Erd\H{o}s and Szemer\'edi asked for a bound exponentially
smaller than \(2^N\) on a sunflower-free family of subsets of
an \(N\)-element set~\cite{ES}. Alon, Shpilka and Umans showed how
both conjectures can be studied using words, and that the fixed-petal
conjecture rules out a group construction proposed by Coppersmith
and Winograd for matrix multiplication~\cite{ASU}. Naslund and
Sawin gave a direct slice-rank proof of the ground-set bound for
three petals~\cite{NS}. Buci\'c, He and Lee obtained
\(3(2-1/(16r^3))^N\) for every \(r\ge3\) by polynomial
interpolation~\cite[Theorem~1]{BHL}. For words of length \(k\) over
\(D\) symbols, their bound has base
\((2r)^{2/(r-1)}D^{1-1/(r-1)}\)~\cite[Theorem~2]{BHL}. The
word estimate proved here has a single base for all alphabets, as
required for the passage to \(k\)-element sets on an arbitrary ground
set~\cite[Theorem~2.6]{ASU}.

